% Part: history
% Chapter: biographies
% Section: rozsa-peter

\documentclass[../../../include/open-logic-section]{subfiles}

\begin{document}

\olfileid{his}{bio}{pet} 

\olsection{روژا پیتر}

\olphoto{peter-rozsa}{روژا پیتر}

روژا پیتر 17 فروری 1905 نوں بڈاپسٹ، ہنگری وچ روژا پولٹسر دے ناں نال جمی۔ اوہ سبھ توں ودھ بازگشتی فنکشناں اُتے اپنے کم لئی جانی جاندی اے۔ ایہ کم بازگشت دا نظریہ اک شعبے وجوں بناؤن لئی ضروری سی۔

پیتر سخت سیاسی حالات وچ وڈی ہوئی۔ اوہدی چڑھدی عمر وچ پہلی عالمی جنگ چل رہی سی۔ پر اوہ بڈاپسٹ دے خوشحال ماریا تیریزیا گرلز سکول وچ پڑھ سکی تے 1922 وچ اوتھوں پڑھائی پوری کیتی۔ فیر اوہنے بڈاپسٹ دی پازمانی پیتر یونیورسٹی (جیہدا ناں بعد وچ لوراند ایوتووش یونیورسٹی رکھیا گیا) وچ پڑھیا۔ اپنے پیو دے زور دین اُتے اوہنے کیمیا پڑھنی شروع کیتی، پر فیر ریاضی ول آ گئی تے 1927 وچ ڈگری لئی۔ بھاویں اوہدے کول ہائی سکول وچ ریاضی پڑھاؤن دی اہلیت سی، اوس ویلے مالی حال بڑا ماڑا سی، کیوں جے وڈی کساد بازاری دا اثر سارے جہان دی معیشت اُتے پیا سی۔ ایس سمیں پیتر نے ریاضی دی نجی استاد وجوں تھوڑے موٹے کم کیتے۔ آخر اوہ ریاضی دی اعلیٰ پڑھائی کرن لئی یونیورسٹی مُڑ آئی۔ اوہدا پہلا ارادہ عدد نظریے وچ کم کرن دا سی، پر جدوں پتا لگیا کہ اوہدے نتیجے پہلاں ای ثابت ہو چکے سن، تاں اوہ ریاضی چھڈن لگی سی۔ اوہنوں گوڈل دے عدم تکمیل دے قضیاں اُتے کم کرن دی ہمت دِتی گئی، تے اوہنے انجانے وچ اوہدے کئی نتیجے وکھرے طریقیاں نال ثابت کر دتے۔ ایس نال اوہدا بھروسا مُڑ آیا، تے ڈیوڈ ہلبرٹ دے ریاضی دیاں بنیاداں والے پروگرام توں تحریک لے کے اوہنے بازگشت دے نظریے اُتے اپنے پہلے مقالے لکھے۔ اوہنے 1935 وچ پی ایچ ڈی لئی، تے 1937 وچ \emph{Journal of
  Symbolic Logic} دی مدیر بنی۔

پیتر دے مُڈھلے مقالے بڑے پیمانے اُتے بازگشتی فنکشناں دے نظریے دی بنیاد رکھن والا کم منّے جاندے نیں۔ \cite{Peter1935a} وچ اوہنے وکھ وکھ قسماں دی بازگشت دا آپسی تعلق ویکھیا۔ \cite{Peter1935b} وچ اوہنے وکھایا کہ بازگشت نال تعریف کیتا اک خاص فنکشن ابتدائی بازگشتی نہیں۔ ایس نال ولہلم آکَرمان دا پہلا نتیجہ سوکھا ہو گیا۔ پیتر دا سوکھا کیتا فنکشن ہُن اکثر آکَرمان فنکشن آکھیا جاندا اے، تے کدی ہور ٹھیک ناں، آکَرمان–پیتر فنکشن، نال وی۔ اوہنے بازگشتی فنکشناں دے نظریے اُتے پہلی کتاب لکھی \citep{Peter1951}۔

اپنے کم دی اہمیت تے اثر دے باوجود پیتر نوں 1945 تک پورے ویلے دی پڑھاؤن والی نوکری نہ ملی۔ دوجی عالمی جنگ وچ ہنگری اُتے نازی قبضے ویلے یہودی مخالف قاعدیاں کرکے پیتر نوں پڑھاؤن دی کھُل نہیں سی۔ 1944 وچ سرکار نے بڈاپسٹ وچ یہودیاں لئی اک گھِری ہوئی بستی بنائی۔ ایہ بستی باقی شہر توں کٹّی ہوئی سی تے اوتھے ہتھیاربند پہرا سی۔ پیتر نوں 1945 وچ بستی آزاد ہون تک اوتھے رہن لئی مجبور کیتا گیا۔ فیر اوہنے بڈاپسٹ ٹیچرز ٹریننگ کالج وچ پڑھایا تے 1955 توں ایوتووش لوراند یونیورسٹی وچ۔ اوہ ریاضی دی اکادمی والی ڈاکٹریٹ لین والی پہلی ہنگیرین عورت ریاضی دان تے ہنگیرین اکادمی آف سائنسز لئی چُنی جان والی پہلی عورت سی۔

پیتر ریاضی دی پرجوش استاد وجوں جانی جاندی سی۔ اوہ بس لیکچر دین دی بجائے اپنے طالب علماں نال ریاضیاتی مسئلیاں دی نوعیت تے سوہپ لبھنا پسند کردی سی۔ ایس لئی اوہدے طالب علم پیار نال اوہنوں «ماسی روزا» آکھدے سن۔ پیتر 1977 وچ~71 سال دی عمر وچ وفات پا گئی۔

\begin{reading}
اوہدی زندگی بارے ہور پڑھن لئی \citep{Oconnor2014} تے \citep{Andrasfai1986} ویکھو۔ \citet{Tamassy1994} نے پیتر نال اک مختصر انٹرویو کیتا۔ ریاضی بارے دلچسپ پڑھائی لئی پیتر دی کتاب \emph{Playing With Infinity} ویکھو \citep{Peter2010}۔
\end{reading}

\end{document}
